\section{自信息 $I(x)$：一条消息值多少 bit}

\begin{idea}{先做一道猜数字的题，你就懂 $I(x)$ 了}
我心里想了一个 $1$ 到 $8$ 之间的整数，你来猜。你每次只能问一个是非题（“是大于 4 吗？”之类），我如实回答。

最聪明的问法是“对半砍”：先问“大于 4 吗？”——无论我怎么答，剩下的可能性都从 8 个减半成 4 个；再砍一刀变 2 个；再砍一刀变 1 个。所以\textbf{最多 3 次一定能猜中}。

现在换个说法：我一开始心想的那个数，在你眼里有 $8$ 种可能，而我一告诉你答案，可能性从 $8$ 变成 $1$。你猜了 $3$ 次，用了 $3$ 个“是/不是”的回答。于是我们很自然地这样规定：
\[
\text{把可能性砍掉一半，就记 }\ 1\ \text{个信息量（1 个二进制单位）},\qquad \text{砍 3 刀} \Rightarrow 3\ \text{个单位}.
\]
因为 $8=2^3$，所以“$8$ 种可能里告诉你哪一个”正好值 $3$ 个单位，而 $3=\lb 8$。
把这件事翻译成公式：\textbf{如果这件事的先验概率是 $p$，那么它发生所提供的信息量是 $-\lb p$}。
\end{idea}

\subsection{自信息的定义}

\begin{defbox}{定义（自信息 / 信息量）}
设某件事 $x$ 发生的概率为 $p(x)$，则它\textbf{发生（被你知道）时所提供的信息量}，称为 $x$ 的\textbf{自信息}，记作
\[
\boxed{\ I(x)=-\lb p(x)=\lb\frac{1}{p(x)}\ }.
\]
单位是 \textbf{bit}（比特）。
\end{defbox}

\begin{idea}{这个式子的三处“人话”翻译}
\begin{enumerate}[leftmargin=1.6em,itemsep=1pt,topsep=2pt]
  \item \textbf{为什么是 $\frac{1}{p}$ 取对数}：$p(x)$ 越小（越意外），$\frac{1}{p(x)}$ 越大，信息量越大——越意外的消息越值钱。概率为 $1$ 的事（“太阳明天会升起来”）谈不上任何信息。
  \item \textbf{为什么外面有个负号}：因为 $p\le1$，所以 $\lb p\le0$，前面加个负号才能让信息量变成\textbf{非负}的数。所以写 $-\lb p$ 和写 $\lb\frac1p$ 完全是一回事，后者更不容易算错。
  \item \textbf{为什么是 log}：这样才具有“\textbf{可加性}”——两条互不相关消息同时告诉你的信息量，等于各自信息量之和（见下面性质 3）。
\end{enumerate}
\end{idea}

\subsection{手把手算几个最基本的自信息}

\begin{calc}{例 1：抛公平硬币}
$p(\text{正面})=\frac12$，于是
\[
I(\text{正面})=-\lb\frac12=\lb 2=1\bit .
\]
\textbf{这就是“1 bit”这个词的由来}：抛一次公平硬币、告诉你结果，你得到的正好是 1 bit 信息。比如硬盘上一个 bit 存的就是“是/否”。
\end{calc}

\begin{calc}{例 2：掷一颗骰子}
$p(\text{出现 }6)=\frac16$，于是
\[
I=-\lb\frac16=\lb 6\approx2.585\bit .
\]
为什么不是整数？因为 $6$ 不是 2 的整数次方，没法用整数个“是/不是”问题问完。$2.585$ 的意思是“平均需要 2.585 个二进制位才能确定结果”。
同理，从 $8$ 匹马中告诉你哪匹赢（等概率）：$-\lb\frac18=\lb8=3\bit$，正好 3 个问题；从 $64$ 个方块中猜中一个：$-\lb\frac1{64}=6\bit$。
\end{calc}

\begin{calc}{例 3：不公平的硬币（这里藏着本课的核心思想）}
设 $p(\text{正})=0.9$，$p(\text{反})=0.1$，则
\[
I(\text{正})=-\lb 0.9\approx0.152\bit,\qquad
I(\text{反})=-\lb 0.1\approx3.322\bit .
\]
\textbf{体会}：如果硬币严重偏向正面，“正面”这个结果几乎不意外，只值 $0.152$ bit；而一旦出现“反面”，那是很惊人的消息，值 $3.322$ bit。
可见\textbf{“意外的消息更值钱”}。那么平均每次抛硬币能收到多少信息？你不能只取一个数——这就是下一讲的“熵”：把 $0.152$ 和 $3.322$ 按 $0.9$ 和 $0.1$ 加权平均。
\end{calc}

\begin{calc}{例 4：天气预报（考试常用数据）}
设明天天气：晴 $p=0.5$、阴 $p=0.25$、雨 $p=0.25$。各事件的自信息为
\begin{align*}
I(\text{晴})&=-\lb0.5=1\bit,\\
I(\text{阴})&=-\lb0.25=-\lb(2^{-2})=2\bit,\\
I(\text{雨})&=-\lb0.25=2\bit .
\end{align*}
（这里用了一个小技巧：$0.25=\frac14=2^{-2}$，所以 $-\lb0.25=2$，不用查表。）
\end{calc}

\subsection{自信息的三个性质（对应“为什么用 log”的常识要求）}

\begin{defbox}{性质 1：单调递减}
概率越大，信息量越小：
\[
p(x_1)<p(x_2)\ \Longrightarrow\ I(x_1)>I(x_2).
\]
人话：越罕见的事，说出来越有信息。
\end{defbox}

\begin{defbox}{性质 2：非负性}
\[
p(x)\le 1\ \Longrightarrow\ I(x)=-\lb p(x)\ge 0,
\]
且 $p(x)=1$ 时 $I(x)=0$。
人话：“早知道了一定会发生的事”不提供任何信息。
\end{defbox}

\begin{defbox}{性质 3：可加性（这条决定了只能用 log）}
若事件 $x$ 与 $y$ 相互独立（$p(x,y)=p(x)p(y)$），则
\[
I(xy)=-\lb p(x,y)=-\lb\bigl[p(x)p(y)\bigr]=-\lb p(x)-\lb p(y)=I(x)+I(y).
\]
验证一下：抛两次公平硬币，两个结果的信息量应为 $1+1=2$ bit；
直接算：$p(\text{正正})=\frac12\cdot\frac12=\frac14$，$I=-\lb\frac14=2\bit$。$\ $对上！
人话：\textbf{两条互不相干的消息，总信息量等于各自之和}——这几乎是任何“信息量”定义都必须满足的要求，而只有对数能做到（若定义 $I=1/p$，则 $\frac{1}{1/4}=4\ne2+2$，不成立）。
\end{defbox}

\subsection{更一般的三种“自信息”}

除了“一个事件的自信息”，书上还定义了另外两种，写法很像，要分清：

\begin{defbox}{联合自信息、条件自信息}
\[
I(x,y)=-\lb p(x,y),\qquad
I(x|y)=-\lb p(x|y),\qquad
I(x;y)=\lb\frac{p(x,y)}{p(x)p(y)} .
\]
\begin{itemize}[leftmargin=1.6em,itemsep=1pt,topsep=1pt]
  \item $I(x,y)$：\textbf{“$x$ 和 $y$ 同时发生”}这个联合事件的信息量；
  \item $I(x|y)$：\textbf{在已知 $y$ 的条件下}，再告诉你“$x$ 发生了”还值多少信息（因为已知 $y$ 后 $p(x|y)$ 可能变大，所以这个值通常比 $I(x)$ 小）；
  \item $I(x;y)$：互信息，衡量“$y$ 帮你在 $x$ 上省了多少”——它\textbf{可正可负}（第5讲讲）。
\end{itemize}
\end{defbox}

\begin{calc}{手把手：用上一讲那张联合分布表算三种自信息}
表：$p(0,0)=\frac13$，$p(0,1)=\frac13$，$p(1,0)=\frac13$，$p(1,1)=0$；边缘 $p(X=0)=\frac23$，$p(X=1)=\frac13$，$p(Y=0)=\frac23$，$p(Y=1)=\frac13$。

（1）$I(X=1)=-\lb\frac13=\lb3\approx1.585$ bit（$\frac13$ 比 $\frac12$ 小，所以比 1 bit 略大）。
（2）联合自信息 $I(X=0,Y=0)=-\lb\frac13=\lb3\approx1.585$ bit。
（3）条件自信息：已知 $Y=0$ 后 $X=0$ 的条件概率是 $\frac12$（上一讲算过），故
\[
I(X=0|Y=0)=-\lb\frac12=1\bit .
\]
对比一下：$I(X=0)=\lb\frac32\approx0.585$ bit，而 $I(X=0|Y=0)=1$ bit。
\textbf{什么？知道 $Y=0$ 之后反而更意外了？}是的，这种情况完全可能：因为 $Y=0$ 这个条件把 $p(X=0)$ 从 $\frac23$ 压到了 $\frac12$。所以\textbf{“条件不增信息量”并不是单条事件成立的}，它只对“平均以后”的熵成立——这是初学者最容易错的地方，第4讲会专门讲。
\end{calc}

\begin{exam}{第2讲：考试怎么用}
\begin{enumerate}[leftmargin=1.6em,itemsep=1pt,topsep=2pt]
  \item \textbf{直接给概率求信息量}：$I=-\lb p$，套公式，注意单位写 bit。答题时写“$I(x)=-\lb 0.25=2$ bit”这样的完整一行即可得分。
  \item \textbf{反着问“需要多少位”}：若信息量是 $I$ bit，则大约需要 $I$ 个二进制位表示（有小数时向上取整）。考试里常问“$n$ 元等概消息需要多少 bit”$\Rightarrow \lb n$。
  \item \textbf{判断题“概率越小信息量越大”}：正确；\textbf{“概率为 0 的事件信息量为无穷大”}：这是数学上的极限（现实中不会发生），一般不出这种题，若问就答“$I$ 随 $p\to0$ 而趋于无穷”。
  \item \textbf{千万别把 $I(x)$ 和 $H(X)$ 混为一谈}：$I(x)$ 是“某一条消息”的信息量（随具体取值变化，有大小不一定相等）；$H(X)$ 是“平均每条”的信息量（一个数）。下一讲就是讲 $H$。
\end{enumerate}
\end{exam}

\begin{qa}{自测：你真的会了吗（有答案）}
\textbf{问}：某信源发出四种符号，概率分别为 $1/2,\ 1/4,\ 1/8,\ 1/8$。求每个符号的自信息，以及（不用算平均）哪个符号最有信息量。

\textbf{答}：$I_1=-\lb\frac12=1$ bit；$I_2=-\lb\frac14=2$ bit；$I_3=I_4=-\lb\frac18=3$ bit。
最有信息量的是概率最小、也就是最罕见的两个符号（各 3 bit）。
顺便说：如果你要设计编码，直觉上也应该把最短的码字给最常出现的符号——这正是第8讲 Huffman 编码的全部思想。
\end{qa}
